\documentclass{article}
\usepackage[utf8]{inputenc}
\usepackage[T1]{fontenc}

\usepackage[english]{babel}
\usepackage{amsmath,amssymb,mathtools}
\usepackage{booktabs}
\usepackage{graphicx}
\usepackage{hyperref}
\usepackage{enumitem}

\graphicspath{{figures/}}

%\pagenumbering{gobble}

\begin{document}


\begin{center}
	{\huge Predicting Wordle Results} \\
		\vspace*{3mm}
		{\large 19M081MMS Mathematical Modeling and Simulations \\
		Palace of Science, Center for Applied Mathematics \\
		\vspace*{3mm}
		17 April 2026. }
\end{center}

\vspace*{10mm}

\begin{tabular}{r l}
\textbf{Students:} &  2025/3228 Mihailo Radović \\
				&	2025/3330 Boško Zlatanović \\
				&	2025/3318 Filip Marčić \\
				&   \\
\textbf{Mentor:} & Prof.~Dr.~Nataša Ćirović
\end{tabular}

\vspace*{5mm}
\textbf{Main responsibilities:}
\vspace*{3mm}

\begin{tabular}{r l}
	Problem presentation:  Boško Zlatanović, Mihailo Radović, Filip Marčić&  \\
	Programming: Boško Zlatanović, Mihailo Radović, Filip Marčić	\\
	Report:	Boško Zlatanović
\end{tabular}

\vspace*{15mm}

\section{Abstract}

This project studies daily Wordle results collected from Twitter posts during
2022. Each daily record contains the solution word, the number of reported
scores, the number of Hard Mode players, and the distribution of outcomes across
one to six tries and failures (denoted $X$). The analysis addresses three
connected questions:

\begin{enumerate}[label=\arabic*)]
  \item how participation evolves over time and what value it is expected to take
        on 1 March 2023,
  \item how the full seven-bin solve distribution of a future solution word can
        be predicted from lexical features alone,
  \item how solution words can be ranked by intrinsic difficulty after removing
        the temporal trend in expected attempts.
\end{enumerate}

The pipeline combines data cleaning, lexical and phonetic feature engineering,
log-transformed ridge regression for participation, multi-output regression
(ridge, softmax, CatBoost) for the solve distribution under Kullback--Leibler
loss, and a detrended quartile-based classifier for difficulty. Final numerical
results, including the forecast for 1 March 2023 and the predicted distribution
for the word \texttt{EERIE}, are reported in the corresponding sections.

\section{Introduction}

\subsection*{Background and motivation}
Wordle is a daily word-guessing game in which players attempt to identify a
five-letter solution word within six tries. During 2022 a large number of
players shared their results on Twitter, producing a public time series of
daily participation and outcome distributions. Modeling this data is of
interest because it combines a clear temporal component (audience growth and
decay) with a per-word lexical component (which words are intrinsically harder
to solve).

\subsection*{Problem formulation}
The problem is not only to describe the observed Wordle data but to build
predictive models for future behavior. The report focuses on:

\begin{itemize}
  \item forecasting the number of reported results on a future date
        (1 March 2023),
  \item predicting the full outcome distribution for a future solution word,
  \item classifying words such as \texttt{EERIE} into a four-level difficulty
        scale,
  \item identifying which lexical and structural characteristics are most
        strongly associated with difficulty.
\end{itemize}

\subsection*{Assumptions}
A key empirical observation, established in the exploratory notebook and used
throughout, is that the reporting population is non-stationary: the share of
reports submitted in Hard Mode rises steadily through 2022 while the absolute
count of Hard Mode players is comparatively flat. This implies that the casual
audience is shrinking and that recent observations are more representative of
the regime to be forecast. We further assume that the seven try-counts can be
treated as a probability vector and that the CMU pronouncing dictionary is a
reasonable, if imperfect, source for phonetic features.

\section{Problem Analysis and Model}

\subsection*{Variables}
The independent variables include the integer day index $t$, the solution word
and its lexical features (rarity, structural, sequence, and phonetic
attributes), and the daily context (number of reported results, Hard Mode
fraction, weekend indicator). The dependent variables are: (i) the number of
reported scores $n_t$ on day $t$, (ii) the seven-bin solve probability vector
$\mathbf{p} = (p_1,\dots,p_6,p_X)$, and (iii) a four-level ordinal difficulty
class. Model parameters include the ridge regularization $\alpha$, CatBoost
hyperparameters, and quartile thresholds derived from the detrended target.

\subsection*{Model}

\paragraph{Data cleaning.}
The raw data set~[1] covers 359 daily contests during 2022. The cleaning step
renames the columns to a consistent lowercase schema, corrects five obvious
transcription errors in the solution words
(\texttt{clen}$\to$\texttt{clean}, \texttt{rprobe}$\to$\texttt{probe},
\texttt{tash}$\to$\texttt{trash}, \texttt{favor\,}$\to$\texttt{favor},
\texttt{marxh}$\to$\texttt{march}), fixes a missing trailing zero in the
\texttt{n\_reported} field for the word \texttt{study} (2569 $\to$ 25690),
removes the single row dated 27 March 2022 whose try percentages do not sum
consistently, and exports the cleaned table as \texttt{data/cleaned.csv}.

Two derived quantities are used repeatedly:
\begin{equation}
  h_t = 100 \cdot \frac{n^{\mathrm{hard}}_t}{n^{\mathrm{reported}}_t},
  \qquad
  t = (\text{date} - \text{date}_{\min}) \ \text{in days}.
\end{equation}

\begin{figure}[h!]
  \centering
  \includegraphics[width=0.55\linewidth]{fig_boxplot_nreported.png}
  \caption{Boxplot of \texttt{n\_reported} over the full 2022 period. A long
  upper tail in the first quarter of 2022 motivates restricting the forecasting
  sample to the post-May regime.}
  \label{fig:boxplot}
\end{figure}

A by-word inspection of the try distribution suggests that, for most words, the
seven bins are reasonably approximated by a discrete bell shape centered between
three and five tries.

\begin{figure}[h!]
  \centering
  \includegraphics[width=0.6\linewidth]{fig_tries_distribution_study.png}
  \caption{Empirical try distribution for the word \texttt{study} overlaid with
  a fitted normal density.}
  \label{fig:tries-normal}
\end{figure}

\paragraph{Model 1: Forecasting participation.}
The target is the number of reported scores $n_t$ on 1 March 2023. The
modeling sample is restricted to dates after 10 May 2022 (235 observations).
A simple linear baseline
\begin{equation}
  n_t = \beta_0 + \beta_1 t + \varepsilon_t
\end{equation}
is replaced by a log-transformed ridge model~[3,6] with $\alpha=12$ on
standardized $t$:
\begin{equation}
  y_t = \log(1+n_t),
  \qquad
  \hat n_t = \exp(\hat y_t) - 1.
\end{equation}
Model selection used a chronological 90/10 train/test split, and the
log-transformed ridge gave a lower test RMSE than the linear baseline.

\begin{figure}[h!]
  \centering
  \includegraphics[width=0.95\linewidth]{fig_nreported_over_time.png}
  \caption{Number of reported results over time. The modeling sample begins on
  10 May 2022.}
  \label{fig:nreported}
\end{figure}

\begin{figure}[h!]
  \centering
  \includegraphics[width=0.95\linewidth]{fig_participation_fit.png}
  \caption{Train/test fit of the exponential ridge model on
  \texttt{n\_reported}.}
  \label{fig:participation-fit}
\end{figure}

\paragraph{Model 2: Predicting the solve distribution.}
The target is the seven-category probability vector
$\mathbf{p} = (p_1, p_2, p_3, p_4, p_5, p_6, p_X)$, with per-row normalization
so that $\sum_k p_k = 1$. The evaluation metric is the Kullback--Leibler
divergence
\begin{equation}
  D_{\mathrm{KL}}(\mathbf{p}\,\|\,\mathbf{q})
  = \sum_{k} p_k \log\!\left(\frac{p_k}{q_k}\right),
\end{equation}
evaluated per row and averaged over a five-fold cross-validation, with both
vectors clipped to a small $\varepsilon = 10^{-9}$.

The feature set is grouped as:
\begin{itemize}
  \item \textbf{Rarity features:} corpus-level word rarity from
        \texttt{wordfreq}~[7] (\texttt{zipf\_frequency}), sum of log letter
        frequencies~[10], inverse-frequency-weighted rarity;
  \item \textbf{Structure features:} number of unique letters, duplicate-letter
        penalty, number of vowels, number of rare letters
        ($\{z,x,q,j,k,v,b,p,y,g\}$);
  \item \textbf{Sequence features:} start- and end-letter scores from empirical
        first- and last-letter frequencies~[10], common-bigram and
        common-trigram counts~[10], positional negative-log-likelihood;
  \item \textbf{Context features:} the total number of reported results for the
        day.
\end{itemize}

Three estimators were compared: per-bin ridge regression~[3,6] with
$\alpha=150$ followed by non-negative clipping and renormalization; a softmax
linear model that passes seven ridge predictions through a row-wise softmax;
and a CatBoost multi-output regressor~[4,5] with \texttt{MultiRMSE} loss tuned
via randomized search. For each estimator a greedy forward feature selection
was run, scored by mean cross-validated KL divergence. If a model returns raw
bin scores $s_1,\dots,s_7$, the prediction is
\begin{equation}
  \hat p_k = \frac{\max(s_k, 0)}{\sum_{j=1}^{7} \max(s_j, 0)}.
\end{equation}

\paragraph{Model 3: Word difficulty classification.}
Each solution word is assigned to one of four ordinal difficulty classes.
Given the empirical distribution $\mathbf{p}_t$, the expected number of
attempts on day $t$ is
\begin{equation}
  \mu_t = \sum_{k=1}^{7} k\,p_{t,k}.
\end{equation}
Because the corpus-wide mean drifts downward over the year, a univariate trend
is removed before classifying:
\begin{equation}
  \mu_t \approx \alpha + \beta\, c_t,
  \qquad
  y^{\mathrm{dtr}}_t = \mu_t - (\hat\alpha + \hat\beta\, c_t),
\end{equation}
where $c_t$ is a normalized contest index in $[0,1]$. The detrended values are
split into quartiles to define the four ordinal classes \textbf{Easy},
\textbf{Medium}, \textbf{Hard}, and \textbf{Very Hard}.

In addition to the features of Model~2, the difficulty model uses per-letter
and per-position frequency statistics~[10] including a positional surprise
score; consonant--vowel pattern rarity; Laplace-smoothed bigram
log-probabilities~[10]; repeated-letter information loss
$\log_2\!\bigl(5!/\prod c_i!\bigr)$, vowel-cluster length, and letter entropy;
optional phoneme-based features~[9] via the \texttt{pronouncing} CMU
wrapper~[8]; and context features (Hard Mode fraction, weekend indicator).
Two regressors are evaluated on $y^{\mathrm{dtr}}$: ridge regression~[3,6] on
standardized features and a random forest~[2,6] with $500$ trees and
\texttt{min\_samples\_leaf}$=5$. The continuous prediction is then mapped back
to one of the four ordinal classes using the quartile thresholds, and agreement
is measured by quadratic-weighted Cohen's $\kappa$.

\subsection*{Computer implementation}
All models are implemented in Python. Data cleaning and exploration use
\texttt{pandas}, \texttt{numpy}, and \texttt{matplotlib}. Ridge regression and
random forest models are taken from \texttt{scikit-learn}~[6]. Gradient
boosting uses \texttt{CatBoost}~[4,5]. Lexical features use
\texttt{wordfreq}~[7] for corpus rarity and \texttt{pronouncing}~[8] for
phonetic attributes. The cleaning notebook produces \texttt{data/cleaned.csv};
the participation, distribution, and difficulty pipelines live in
\texttt{1.ipynb}, \texttt{2.ipynb}, and \texttt{3.py} respectively. The main
practical difficulties encountered were missing phonetic entries in the CMU
dictionary for invented Wordle words and the need to enforce a valid
probability vector at prediction time.

\subsection*{Simulation and results}

Refit on the full filtered sample, the participation model predicts
\begin{equation}
  \widehat{n}_{2023\text{-}03\text{-}01} \approx 15{,}607
  \text{ reported results.}
\end{equation}

The best forward-selected subset for the per-bin ridge baseline of Model~2,
\(\{\text{n\_unique\_letters},\ \text{n\_rare\_letters},\ \text{rarity},\)
\(\text{positional\_entropy},\ \text{start\_letter\_score}\}\), achieved a
mean KL of $0.0577$. The softmax linear model gave a higher KL
(mean $\approx 0.365$). The CatBoost multi-output regressor produced the
lowest KL among the three approaches and was used for the final
\texttt{EERIE} prediction.

The word \texttt{EERIE} contains three vowels, the repeated letter \texttt{e},
and an unusually low number of unique letters (three). The CatBoost prediction
for 1 March 2023 is reported in Table~\ref{tab:eerie}.

\begin{table}[h!]
  \centering
  \begin{tabular}{lc}
    \toprule
    Bin & Predicted probability \\
    \midrule
    1 try            & 0\,\% \\
    2 tries          & 5\,\% \\
    3 tries          & 18\,\% \\
    4 tries          & 32\,\% \\
    5 tries          & 29\,\% \\
    6 tries          & 13\,\% \\
    7 or more ($X$)  & 2\,\% \\
    \bottomrule
  \end{tabular}
  \caption{Predicted solve distribution for \texttt{EERIE} using the
  forward-selected CatBoost model.}
  \label{tab:eerie}
\end{table}

The implied expected number of attempts is
\(\sum_k k\,\hat p_k \approx 4.30\), placing \texttt{EERIE} above the
contemporaneous mean.

\begin{figure}[h!]
  \centering
  \includegraphics[width=0.6\linewidth]{fig_eerie_distribution.png}
  \caption{Predicted solve distribution for \texttt{EERIE} (CatBoost).}
  \label{fig:eerie-dist}
\end{figure}

Under the detrended scheme of Model~3, \texttt{EERIE} is predicted in the upper
tail of the difficulty scale by both Ridge and Random Forest, consistent with
its three unique letters, two vowel clusters, and the repeated initial
\texttt{ee} bigram. The classifier therefore places it in the \textbf{Hard} or
\textbf{Very Hard} bracket.

\begin{figure}[h!]
  \centering
  \includegraphics[width=0.85\linewidth]{fig_ridge_coefficients.png}
  \caption{Ridge coefficients on standardized features (positive = harder).}
  \label{fig:ridge-coef}
\end{figure}

\begin{figure}[h!]
  \centering
  \includegraphics[width=0.85\linewidth]{fig_rf_importance.png}
  \caption{Top random-forest feature importances.}
  \label{fig:rf-importance}
\end{figure}

\subsection*{Analysis and discussion}

The three models are complementary: the participation model explains how many
people are expected to report; the distribution model predicts how those
reports spread across outcomes; the difficulty model explains \emph{why}
certain solution words shift the distribution toward worse outcomes. The signed
Ridge coefficients and the random-forest importances agree on the qualitative
picture: difficulty is not driven by a single variable. Repeated letters, low
mean letter frequency, unusual consonant--vowel patterns, and weak bigram
structure all contribute positively to the predicted detrended score.
Conversely, words with high letter entropy and a common opening letter score
as easier.

The main modeling lessons are:
\begin{enumerate}
  \item Temporal regime matters: restricting the participation sample to the
        post-May 2022 period materially improved the 2023 forecast.
  \item Word structure matters more than raw letter rarity in isolation;
        repetition and positional surprise carry the most signal.
  \item Predicting the full seven-bin distribution is more informative than
        predicting a single mean attempt count, because the failure tail $p_X$
        carries information about audience self-selection that the mean
        absorbs.
\end{enumerate}

Several limitations should be acknowledged: the data set spans a single year,
so the model cannot distinguish a long-term decline from a one-off
post-pandemic correction; phonetic features rely on the CMU dictionary, which
fails for many invented or proper-noun Wordle words; and the difficulty target
is constructed from the same outcome distribution that Model~2 predicts, so
Models~2 and~3 share information and cross-comparison should be interpreted
with care.

\section{Conclusion}

Wordle results admit accurate prediction from a small set of carefully chosen
features. A log-transformed ridge regression provides a robust participation
forecast (\(\widehat{n}_{2023\text{-}03\text{-}01}\approx 15{,}607\)). Lexical
features combined with multi-output regression deliver a calibrated solve
distribution (best mean KL $\approx 0.058$ for the per-bin Ridge baseline,
lower for CatBoost). A detrended regression converted to quartile classes
gives a transparent difficulty ranking. For \texttt{EERIE} all three pieces
agree: a low-volume, hard-to-solve word with predicted distribution mass
centered on four to five tries.

The strongest conclusion is that both the calendar date and the lexical form
of the solution word influence player behavior, so a combined time-series and
feature-based modeling approach is more appropriate than a purely static one.
The model can be verified by holding out later weeks of 2022 and comparing
predicted versus observed participation and distributions. Sensitivity to the
assumed regime change date (10 May 2022) and to the choice of feature subset
is moderate but non-negligible. Future improvements could include nonlinear
time trends, uncertainty intervals for the participation forecast, and richer
linguistic features such as syllable structure or context-aware phoneme
statistics.

\begin{thebibliography}{99}
  \bibitem{comap} COMAP, Inc. \href{https://www.contest.comap.com/undergraduate/contests/mcm/contests/2023/problems/}{\emph{MCM/ICM 2023 Problem~C: Predicting Wordle Results}}
        (data file \texttt{Problem\_C\_Data\_Wordle.xlsx}; daily results
        7~January--31~December 2022).

  \bibitem{breiman} Breiman, L. (2001). \href{https://doi.org/10.1023/A:1010933404324}{Random Forests}.
        \emph{Machine Learning}, 45(1), 5--32.

  \bibitem{esl} Hastie, T., Tibshirani, R., and Friedman, J. (2009).
        \href{https://doi.org/10.1007/978-0-387-84858-7}{\emph{The Elements of
        Statistical Learning: Data Mining, Inference, and Prediction}}
        (2nd ed.). Springer.

  \bibitem{catboost} Prokhorenkova, L., Gusev, G., Vorobev, A., Dorogush,
        A.~V., and Gulin, A. (2018). \href{https://arxiv.org/abs/1706.09516}{CatBoost:
        Unbiased Boosting with Categorical Features}. \emph{Advances in Neural
        Information Processing Systems (NeurIPS)}~31.

  \bibitem{catboostdocs} \href{https://catboost.ai/}{CatBoost documentation}.

  \bibitem{sklearn} Pedregosa, F., et al. (2011).
        \href{https://jmlr.org/papers/v12/pedregosa11a.html}{Scikit-learn:
        Machine Learning in Python}. \emph{Journal of Machine Learning
        Research}, 12, 2825--2830.

  \bibitem{wordfreq} Speer, R. \href{https://github.com/rspeer/wordfreq}{\texttt{wordfreq}:
        Word-frequency data for many languages}.

  \bibitem{cmudict} Carnegie Mellon University.
        \href{http://www.speech.cs.cmu.edu/cgi-bin/cmudict}{\emph{The CMU
        Pronouncing Dictionary}}, accessed via the
        \href{https://github.com/aparrish/pronouncing}{\texttt{pronouncing}}
        Python package (A.~Parrish).

  \bibitem{phonemes} DSF Literacy Resources.
        \href{https://www.readingrockets.org/sites/default/files/migrated/the-44-phonemes-of-english.pdf}{\emph{The
        44 Sounds (Phonemes) of English}}. Reading Rockets.

  \bibitem{letterfreq} Oxford College Math Center, Emory University.
        \href{https://mathcenter.oxford.emory.edu/site/math125/englishLetterFreqs/}{\emph{Letter
        Frequencies in English}}.
\end{thebibliography}

\appendix

\section{Appendix}

The source code is provided as an attachment to this report and does not need
to be included in the same file. Notebooks \texttt{1.ipynb} and
\texttt{2.ipynb} contain the participation and distribution pipelines; the
script \texttt{3.py} contains the difficulty pipeline. The cleaned data set is
exported as \texttt{data/cleaned.csv}.

\vspace{0.5em}

\noindent\textbf{GitHub Repository:}
\href{https://github.com/mradovic38/comap-2023-mcm-wordle}
{github.com/mradovic38/comap-2023-mcm-wordle}


\end{document}
